\section{Introduction}\label{sec:introduction}

For a countable discrete group \(G\), let \(C_r^*(G)\) be the norm
closure of its complex group ring in the left regular representation.
The reduced Baum--Connes assembly map is
\[
\mu_j^G:K_j^G(\underline EG)\longrightarrow K_j(C_r^*(G)),
\qquad j=0,1.
\]
Here \(\underline EG\) is the classifying space for proper actions:
its fixed-point sets are contractible for finite subgroups and empty for
infinite subgroups. The domain is equivariant \(K\)-homology with
\(G\)-compact supports, and assembly takes geometric classes to their
analytic indices. The conjecture asserts that \(\mu_j^G\) is an
isomorphism for \(j=0,1\). We use integral complex \(K\)-theory, with
coefficient algebra \(\C\) carrying the trivial action; no
finite-dimensional or cocompact model for \(\underline EG\) is assumed.

Baum and Connes first formulated the conjecture in a manuscript
circulated in 1982 and published in 2000 \cite{BaumConnes2000}.
Their geometric formulation connects elliptic indices with the
\(K\)-theory of reduced crossed products. The formulation of
Baum--Connes--Higson uses the universal proper space and equivariant
Kasparov theory to state this index problem for second-countable locally
compact groups \cite[Conjecture~3.15]{BCH1994}.
The stronger formulation allows a separable \(G\)-\(C^*\)-algebra \(A\)
as coefficient, with target \(K_j(A\rtimes_rG)\)
\cite[Conjecture~9.6]{BCH1994}.

For the discrete groups considered here, important positive results
cover both flexible Hilbert-space geometry and hyperbolic geometry.
Higson--Kasparov prove assembly with coefficients for groups acting
properly by affine isometries on Hilbert space, including countable
amenable groups \cite[Theorem~9.1 and Corollary~9.2]{HigsonKasparov2001}.
Using Lafforgue's Banach \(KK\)-theory, Mineyev--Yu prove the
coefficient-free conjecture for hyperbolic groups and their subgroups
\cite[Theorems~19 and~20]{MineyevYu2002}; Lafforgue later proves
the conjecture with coefficients for hyperbolic groups
\cite[Theorem~0.4]{Lafforgue2012}.

The distinction between trivial and general coefficients is central to
this history.
Higson--Lafforgue--Skandalis use expanders and spectral projections to
detect failures of exactness in reduced crossed products, obtaining
counterexamples for groupoids and for discrete groups with suitable
commutative coefficients \cite[Section~7 and Remark~12]{HLS2002}.
Those constructions do not give a counterexample to the coefficient-free
group assertion above.

Our obstruction uses the canonical trace
\(\tau_G(\sum_g a_g g)=a_1\), normalized by \(\tau_G(1)=1\).
L\"uck's trace theorem implies that its value on every degree-zero
assembly class is rational. Thus a projection in \(C_r^*(G)\) with
irrational trace would obstruct surjectivity. We construct such a
projection.

\begin{theorem}\label{thm:main}
There exist a finitely generated discrete group \(\Gsur\) and a projection
\(b_0\in C_r^*(\Gsur)\) with irrational canonical trace. Moreover,
\[
[b_0]\notin\mu_0^{\Gsur}\bigl(K_0^{\Gsur}(\underline E\Gsur)\bigr).
\]
In particular, the coefficient-free reduced Baum--Connes conjecture for
countable discrete groups is false.
\end{theorem}

The trace restriction just used holds without a surjectivity assumption.
The extension of \(\tau_G\) to matrices is the sum of the diagonal
traces, and \(\tau_{G*}\) denotes the induced homomorphism on \(K_0\).

\begin{theorem}[L\"uck]\label{thm:trace-image}
For every countable discrete group \(G\),
\[
\tau_{G*}\!\left(\mu_0^G(K_0^G(\underline EG))\right)\subseteq\Q.
\]
\end{theorem}

\begin{proof}[Deduction from the trace theorem]
Put
\(\Lambda_G=\Z[1/|F|:F\le G\text{ finite}]\subseteq\Q\).
L\"uck's Theorem~0.3 \cite{Luck2002} states that the trace image of assembly
after extension of scalars to \(\Lambda_G\) is \(\Lambda_G\).
On \(1\otimes x\), this composite has value
\(\tau_{G*}(\mu_0^G(x))\). Hence that value belongs to
\(\Lambda_G\subseteq\Q\) for every integral class \(x\), without
assuming surjectivity of assembly or injectivity of localization.

For precision, our domain is the direct limit of
\(KK_j^G(C_0(Z),\C)\) over closed invariant \(G\)-compact subspaces
of a proper \(G\)-CW model. Assembly is reduced Kasparov descent
followed by the product with a cutoff projection and this direct limit.
L\"uck uses the same proper equivariant \(KK\)-homology and analytic
index \cite[Sections~4 and~5]{Luck2002}. His cocompact subcomplexes
are cofinal among these supports: the quotient is a CW complex, and
its compact subsets lie in finite subcomplexes. Comparisons with the
Davis--L\"uck formulation are given in
\cite[Corollary~8.4]{HambletonPedersen2004} and
\cite[Theorem~5.3]{Kranz2021}.
\end{proof}

\paragraph{Kernel dimensions and the spectral gap.}
The distinction between an individual kernel and a \(K\)-theory index
already appears in Atiyah's question about irrational
\(L^2\)-Betti numbers of coverings
\cite[Section~6.4, question~(iii)]{Atiyah1976}.
Grigorchuk and \.{Z}uk's lamplighter spectral calculation led
Grigorchuk--Linnell--Schick--\.{Z}uk to examples with kernel dimension
\(1/3\), although every finite subgroup had order a power of two
\cite{GLSZ2000}. Dicks--Schick gave an elementary spectral decomposition
using finite path operators and projections \cite{DicksSchick2002}.
Austin then obtained irrational kernel dimensions by varying
translation-invariant quotients of finite-field lamp modules
\cite[Theorem~1.1 and Corollary~1.2]{Austin2013}.
Fourier transformation turns these operators into local rules on the
compact dual of the module; this viewpoint also underlies the present
construction. Pichot--Schick--\.{Z}uk refined the method to compute
explicit transcendental values \cite{PichotSchickZuk2015}.

These results concern projections in group von Neumann algebras, which
need not belong to reduced group \(C^*\)-algebras. The difference is
particularly clear in Grabowski's construction of integral group-ring
matrices with transcendental kernel dimension over the ordinary
lamplighter groups \((\Z/m\Z)\wr\Z\), for integers \(m>1\)
\cite[Theorem~2]{Grabowski2016}. Those groups are amenable and satisfy
Baum--Connes. By Theorem~\ref{thm:trace-image}, their irrational kernel
projections cannot lie in matrix algebras over their reduced group
\(C^*\)-algebras.

A uniform spectral gap at zero supplies the required \(C^*\)-algebra
property: the indicator of zero becomes continuous on the spectrum.
Li--Nowak--Pooya use this mechanism for cohomological Laplacians and
higher Kazhdan projections, relating gaps and trace restrictions to
surjectivity of assembly \cite[Propositions~1,~5 and~7]{LiNowakPooya2024}.
Here the operator comes from local finite-field rules. Its kernel
consists of constants on designated expander components, and its trace
is computed from their pattern probabilities. Beyond the gaps on those
expanders, the construction must provide a uniform positive lower bound
on the entire remaining graph.

We obtain that bound using labelled expanders as graphical relators,
following the strategy of Gromov \cite{Gromov2003}. Osajda's stronger
labelling theorem provides one finite alphabet, globally distinguishable
long paths, and isometrically embedded relators
\cite[Theorems~2.7 and~3.2]{Osajda2020}. To pass from separate
expander estimates to the full operator, we prove a geometric bound on
simultaneous intersections of translated relators and use it in a
global cycle-rank estimate.

\paragraph{Construction and proof.}
The two ingredients that distinguish the construction are exact local
module relations and a global spectral estimate. Both use the same
geometric ordering, but neither follows from pairwise small overlaps
alone.
We first choose finite \(d\)-regular expanders \(\Theta_i\) with large girth and
label them by Osajda's theorem so that they embed isometrically in a Cayley
graph \(X=\Cay(\Gamma,S)\) of a finitely generated group \(\Gamma\). A \emph{copy} is any left translate of an embedded
\(\Theta_i\). Coning off every copy gives a uniformly hyperbolic graph.
An avoidance estimate then shows that, relative to any root, all
intersections of a copy with other copies whose centers are no farther from
the root lie in one small intrinsic tree ball.

This simultaneous intersection bound has two uses. It makes the finite
cycle spaces of copies an internal direct sum, and it prevents new local
linear relations when codes on copies are imposed on a lamp module.
For one prime \(p\), codes over \(\F_p\) of large dual distance give
exact finite-pattern distributions;
large primal distance forces a positive proportion of discrepancies on
every copy whose pattern is not the designated one.

The resulting group is \(\Gsur=V\rtimes\Gamma\), with \(V\) a finite-field
module. On its compact dual, a coordinate value specifies allowed
generator labels. The operator \(T\in C_r^*(\Gsur)\) is \(d\) minus the
adjacency operator of the enabled graph. Copies carrying the designated
pattern are isolated expander components. To control everything else, we
eliminate copy coordinates in finite cycle spaces and charge their ranks
to disjoint exterior vertex sets. This yields a uniform spectral gap at
zero and hence \(b_0=\one_{\{0\}}(T)\in C_r^*(\Gsur)\) by continuous
functional calculus. Its trace is a series \(\sum_i p^{-k_i}\) whose
base-\(p\) digits are not eventually periodic; here \(k_i\) are the
code dimensions, selected to have increasingly large gaps.

The external geometric inputs are recalled where used. The diagram,
module, and spectral arguments are given in full, including the
simultaneous overlap bound needed to pass from individual copies to the
whole enabled graph.

\paragraph{Organization.}
Section~\ref{sec:geometry} proves the coned geometry and simultaneous
overlap estimate. Section~\ref{sec:data} chooses the labelled expanders
and derives the finite cycle-space decomposition.
Section~\ref{sec:module} constructs the codes, the module, and its exact
finite marginals. Section~\ref{sec:spectrum} proves the spectral gap,
with Figure~\ref{fig:cycle-charging} illustrating the cycle-rank argument.
Section~\ref{sec:trace} computes the irrational trace and applies
Theorem~\ref{thm:trace-image}.
